% -----------------------------------------------------------------------------
% SECTION: SCOPE & INTRODUCTION (ECSS-E-ST-10-06C ANNEX A SECTION <1>)
% -----------------------------------------------------------------------------
\section{Scope}

This document (\texttt{EWAS-HSB-SSEG-TRS-03}) is the Level-2 \textbf{Technical Requirements Specification (\abbr{TRS})} governing Configuration Item \textbf{\texttt{CI-03}: Space Segment (\abbr{SSEG}, \abbr{WBS} Element \texttt{123456.03})} of the \textbf{Early Warning Alarm System (\abbr{EWAS})}. Prepared in strict conformance with \textbf{ECSS-E-ST-10C} (\textbf{AD-02}) and \textbf{ECSS-E-ST-10-06C Annex A} (\textbf{AD-03}), this specification establishes the complete, frozen set of technical requirements necessary to design, manufacture, integrate, verify, launch, commission, operate, and dispose of the \abbr{EWAS} Space Segment as an integrated system element.

The scope of this \abbr{TRS} encompasses all space-borne elements of the two-segment \abbr{EWAS} architecture, specifically:
\begin{enumerate}
    \item \textbf{Hybrid Orbital Constellation Architecture:} The coordinated astrodynamic configuration, orbital slot allocation, multi-plane phasing, and coverage availability of the dual-tier constellation comprising \textbf{six (6) Geostationary Earth Orbit (\abbr{GEO}) detection satellites (\texttt{SATG})} operating at 35,786~km altitude (spaced uniformly at $60^\circ$ longitude intervals) and \textbf{seventy-two (72) Low Earth Orbit (\abbr{LEO}) tracking satellites (\texttt{SATL})} deployed at 1,000~km circular altitude in a high-inclination ($78^\circ\text{--}82^\circ$) Walker Delta $72/8/1$ pattern (8 orbital planes with 9 satellites per plane), satisfying top-level constellation requirements \texttt{R-005} through \texttt{R-008} of \textbf{AD-01}.
    \item \textbf{Space Segment Mission Chain and Latency Allocations:} The functional and performance allocations governing boost-phase Short-Wave/Medium-Wave Infrared (\abbr{SWIR}/\abbr{MWIR}) plume detection, real-time onboard proton-despiking and false-alarm filtering (\texttt{R-009}, \texttt{R-010}, \texttt{R-019}--\texttt{R-022}), immediate candidate-event reporting from \abbr{GEO} to the Ground Segment (\texttt{R-042}), reception of Ground Segment tracking commands (\texttt{R-043}, \texttt{R-044}), agile payload optical Line-of-Sight (\abbr{LOS}) acquisition within $< 5\text{ s}$ and pointing stabilization $< 2\text{ arcsec/s}$ by two selected \abbr{LEO} spacecraft (\texttt{R-024}), synchronized stereoscopic tracking, inter-satellite exchange of timestamped \abbr{LOS} vectors to compute an onboard preliminary 3D trajectory estimate (\texttt{R-023}), and encrypted downlink of both the timestamped stereoscopic \abbr{LOS} observation packages (with covariance metadata) and onboard preliminary trajectory estimate to the Ground Segment Data Processing System for authoritative multi-sensor fusion (\texttt{R-012}, \texttt{R-026}, \texttt{R-027}, \texttt{R-045}) within the Space Segment's allocated share of the $\le 60\text{ s}$ end-to-end warning latency budget (\texttt{R-011}).
    \item \textbf{Mandatory Space Segment Inter-Satellite Link (\abbr{ISL}) Network:} The optical and radio-frequency (\abbr{OISL}/\abbr{RF} \abbr{ISL}) mesh across the 72 \abbr{LEO} nodes and the cross-layer \abbr{GEO}--\abbr{LEO} relay links connecting \texttt{SATL} and \texttt{SATG}, together with cryptographic security (\texttt{R-031}), anti-jamming resilience (\texttt{R-032}), and near-real-time telemetry/command transport (\texttt{R-033}, \texttt{R-044}).
    \item \textbf{Standardized Platform Envelopes, Environment, and Sustainability:} Common platform and subsystem boundaries under the project's \textbf{``Careful \abbr{COTS}'' / Radiation-Tolerant \abbr{COTS}} philosophy, explicitly covering the \emph{Structures and Mechanisms Subsystem (\abbr{STR/MEC})} (including stray-light baffles per \texttt{R-025}, micro-vibration isolation, rideshare dispenser interfaces, and localized \textrm{Al}/\textrm{Ta} radiation spot shielding), \emph{Thermal Control Subsystem (\abbr{TCS})}, \emph{Attitude Determination and Control Subsystem (\abbr{ADCS})}, \emph{Electrical Power Subsystem (\abbr{EPS})}, \emph{Propulsion Subsystem (\abbr{PROP})}, and the architecturally separated \emph{On-Board Data Handling (\abbr{OBDH})} and \emph{Communications (\abbr{TT\&C}/\abbr{ISL})} subsystems, designed for operational lifetimes of 7~years (\abbr{LEO}) and 15~years (\abbr{GEO}) (\texttt{R-016}) and mandatory active end-of-life disposal (\texttt{R-018}, \textbf{AD-12}).
\end{enumerate}

\subsection{Position in the EWAS Specification Tree and Hierarchy}

Per \textbf{ECSS-E-ST-10C} Clause 5.2 (\textbf{AD-02}), the \abbr{EWAS} technical documentation is organized into a hierarchical specification tree aligned with the project Work Breakdown Structure (\abbr{WBS}, Chapter 10 of \textbf{AD-01}):
\begin{itemize}
    \item \textbf{Upper-Tier Parent Baseline (Level 1 --- System Level):} This \abbr{TRS} flows down directly from the \emph{EWAS Overall System Requirements} (\texttt{R-001} through \texttt{R-049}, Chapter 5 of \textbf{AD-01}), the \emph{Concept of Operations} (Chapter 3 of \textbf{AD-01}), and the \emph{Functional Analysis} (Chapter 6 of \textbf{AD-01}). Every technical requirement specified in Section 7 of this document maintains explicit backward traceability to one or more system-level requirements in \textbf{AD-01}.
    \item \textbf{Peer Segment Specifications (Level 2 --- Segment Level):} This document operates at the same hierarchical tier as the \emph{Ground Segment Technical Requirements Specification} (\abbr{WBS} Element \texttt{123456.02}) and is complemented by the \emph{Space Segment Design Definition File} (\texttt{EWAS-HSB-SSEG-DDF-03}), the \emph{Space Segment Design Justification File} (\texttt{EWAS-HSB-SSEG-DDJF-03}), and the segment-level \emph{Interface Control Documents} (\abbr{WBS} Elements \texttt{123456.05.02.01} through \texttt{123456.05.02.04}).
    \item \textbf{Subordinate Lower-Tier Specifications (Level 3 --- Spacecraft \& Subsystem Level):} This \abbr{TRS} defines requirements at the integrated Space Segment boundary (\texttt{123456.03 SSEG}) and establishes the top-level allocation envelopes for its two subordinate flight system Configuration Items: \textbf{Satellite \abbr{LEO} (\texttt{SATL}, \abbr{WBS} \texttt{123456.03.01})} and \textbf{Satellite \abbr{GEO} (\texttt{SATG}, \abbr{WBS} \texttt{123456.03.02})}. Detailed equipment-level and component-level procurement specifications for individual \texttt{SATL} and \texttt{SATG} hardware/software units are allocated to lower-tier subsystem specifications and are not duplicated herein unless required to guarantee cross-layer interoperability, constellation-level performance, or common environmental/product-assurance compliance.
\end{itemize}

\begin{table}[H]
\centering
\caption{EWAS Space Segment (\texttt{CI-03}) WBS Hierarchy and Specification Scope Boundary}
\label{tab:sseg_wbs_scope}
{\renewcommand{\arraystretch}{1.22}%
\arrayrulecolor{TableBorderBlue}%
\rowcolors{2}{white}{TableRowLightBlue}%
\noindent
\begin{tabularx}{\textwidth}{|>{\raggedright\arraybackslash}p{2.5cm}|>{\raggedright\arraybackslash}p{2.7cm}|>{\raggedright\arraybackslash}p{3.9cm}|>{\raggedright\arraybackslash}X|}
    \hline
    \rowcolor{TableHeaderBlue}
    \thead{WBS Code} & \thead{Level / CI} & \thead{Element Title} & \thead{Allocation within \texttt{EWAS-HSB-SSEG-TRS-03}} \\ \hline
    \textbf{123456.00} & Level 1 & EWAS System & Parent system baseline (\texttt{R-001}--\texttt{R-049}, \textbf{AD-01}); source of top-level requirement flow-down. \\ \hline
    \textbf{123456.02} & Level 2 & Ground Systems & External peer segment; interacts with \abbr{SSEG} across the Space-to-Ground \abbr{TT\&C} and mission data interface. \\ \hline
    \textbf{123456.03} & Level 2 (\texttt{CI-03}) & \textbf{Space Segment (\abbr{SSEG})} & \textbf{Primary subject of this TRS:} hybrid constellation geometry, cross-layer \abbr{ISL}, latency budgets, and segment envelopes. \\ \hline
    \textbf{123456.03.01} & Level 3 (\texttt{SATL}) & Satellite \abbr{LEO} (72~nodes) & Subordinate flight segment at 1,000~km (Walker Delta $72/8/1$); platform and \abbr{NFOV} tracking payload envelopes defined herein. \\ \hline
    \textbf{123456.03.02} & Level 3 (\texttt{SATG}) & Satellite \abbr{GEO} (6~nodes) & Subordinate flight segment at 35,786~km ($60^\circ$ spacing); platform and \abbr{WFOV} detection payload envelopes defined herein. \\ \hline
\end{tabularx}}
\end{table}

\newpage

% -----------------------------------------------------------------------------
% INTRODUCTION (ECSS-E-ST-10-06C ANNEX A SECTION <1>: INTRODUCTION)
% -----------------------------------------------------------------------------
\section{Introduction}

\subsection{Purpose and Background}

The \textbf{Early Warning Alarm System (\abbr{EWAS})} is an interdisciplinary space systems engineering project developed jointly at \textbf{Hochschule Bremen (HSB)} and \textbf{Gdańsk University of Technology (Politechnika Gdańska --- PG)}. The mission is motivated by the acute security challenges facing the Baltic Sea region ($53^\circ\text{N}\text{--}66^\circ\text{N}$, $10^\circ\text{E}\text{--}32^\circ\text{E}$), where short flight times of short- and medium-range ballistic missiles (as well as longer-range threats originating outside the region, \texttt{R-049}) leave terrestrial radar networks with severely compressed reaction windows due to Earth-curvature horizon masking and terrain obstruction.

To overcome these physical limitations, \abbr{EWAS} establishes a sovereign, high-availability ($> 99.9\%$, \texttt{R-015}) space-based infrared observation and tracking architecture dedicated strictly to \textbf{civil protection, national crisis management, and emergency preparedness}. In accordance with \texttt{R-001} and \texttt{R-002}, \abbr{EWAS} is strictly a non-offensive early warning system: it does not provide weapon-guidance data, fire-control solutions, or interceptor command-and-control, nor is it required to perform final warhead yield characterization or post-impact damage assessment. Its core operational mandate is to detect ballistic missile launches during the high-radiance boost phase from Geostationary Earth Orbit (\abbr{GEO}), cue two suitable Low Earth Orbit (\abbr{LEO}) satellites via direct ground links or a \textbf{mandatory Space Segment Inter-Satellite Link (\abbr{ISL}) network} (\texttt{R-043}, \texttt{R-044}), execute rapid ($< 5\text{ s}$) optical Line-of-Sight (\abbr{LOS}) acquisition and synchronized stereoscopic tracking (\texttt{R-024}), compute an onboard preliminary 3D trajectory estimate via inter-satellite \abbr{LOS} vector exchange (\texttt{R-023}), and downlink both the timestamped stereoscopic \abbr{LOS} observation packages and preliminary trajectory estimate to the Ground Segment Data Processing System for authoritative multi-sensor fusion (\texttt{R-012}, \texttt{R-026}, \texttt{R-027}). Immediately following this authoritative two-\abbr{LEO} ground assessment, the Ground Segment Early Warning System automatically dispatches a preliminary warning to authorized national authority interfaces within a total end-to-end latency of $\le 60\text{ seconds}$ from initial \abbr{GEO} detection (\texttt{R-011}, \texttt{R-045}).

In compliance with Clause \texttt{<1>a} of \textbf{ECSS-E-ST-10-06C Annex A} (\textbf{AD-03}), the purpose and objective of this Technical Requirements Specification are:
\begin{itemize}
    \item To translate the system-level mission needs, operational scenarios, and requirements (\texttt{R-001} through \texttt{R-049} in \textbf{AD-01}) into unambiguous, quantifiable, attainable, and verifiable technical requirements for the \abbr{EWAS} Space Segment (\texttt{CI-03}, \abbr{WBS} \texttt{123456.03});
    \item To establish the technical reference baseline for the qualification of the Space Segment design (\texttt{EWAS-HSB-SSEG-DDF-03} and \texttt{EWAS-HSB-SSEG-DDJF-03}) and for formal verification and acceptance testing (\textbf{AD-04}, \textbf{AD-05}); and
    \item To define the permissible functional, physical, environmental, operational, and sustainability boundaries governing the 6 \abbr{GEO} (\texttt{SATG}) and 72 \abbr{LEO} (\texttt{SATL}) spacecraft throughout their complete life profile, from on-ground cleanroom integration through 7-year (\abbr{LEO}) and 15-year (\abbr{GEO}) nominal operations to active end-of-life disposal (\texttt{R-016}, \texttt{R-018}).
\end{itemize}

\subsection{Document Structure and Tailoring to ECSS-E-ST-10-06C Annex A}

This specification preserves the standardized HSB/PG project document layout while implementing every normative clause of the \textbf{ECSS-E-ST-10-06C Annex A Technical Requirements Specification Document Requirements Definition (DRD)} (\textbf{AD-03}). Table~\ref{tab:ecss_drd_mapping} defines the exact correspondence between the sections of this document and the mandatory clauses \texttt{<1>} through \texttt{<7>} of \textbf{ECSS-E-ST-10-06C Annex A}.

\begin{table}[H]
\centering
\caption{Mapping between HSB/PG Template Sections and ECSS-E-ST-10-06C Annex A DRD Clauses}
\label{tab:ecss_drd_mapping}
{\renewcommand{\arraystretch}{1.22}%
\arrayrulecolor{TableBorderBlue}%
\rowcolors{2}{white}{TableRowLightBlue}%
\noindent
\begin{tabularx}{\textwidth}{|>{\raggedright\arraybackslash}p{3.5cm}|>{\raggedright\arraybackslash}p{4.2cm}|>{\raggedright\arraybackslash}X|}
    \hline
    \rowcolor{TableHeaderBlue}
    \thead{ECSS-E-ST-10-06C} & \thead{HSB/PG TRS Section} & \thead{Normative Content Summary} \\ \hline
    \textbf{Clause \texttt{<1>}} \newline Introductions & \textbf{Scope} \& \textbf{Introduction} & Purpose, objective, specification tree hierarchy (\texttt{CI-03}), ECSS DRD mapping, and normative requirement conventions. \\ \hline
    \textbf{Clause \texttt{<2>}} \newline Applicable \& Ref.\ Docs & \textbf{References} (\abbr{AD}, \abbr{RD}, \& Abbreviations) & Normative Applicable Documents (\texttt{AD-01}--\texttt{AD-14}), informative Reference Documents (\texttt{RD-01}--\texttt{RD-06}), and abbreviations. \\ \hline
    \textbf{Clause \texttt{<3>}} \newline User's need presentation & \textbf{Main Content --- Sec.\ 3} \newline User's Need Presentation & Baltic security context, non-offensive boundaries, role of \abbr{SSEG} in the two-segment architecture, and external interfaces. \\ \hline
    \textbf{Clause \texttt{<4>}} \newline Selected concept / product & \textbf{Main Content --- Sec.\ 4} \newline Selected Concept Presentation & 6~\abbr{GEO} + 72~\abbr{LEO} (1,000~km Walker Delta) hybrid architecture, mandatory \abbr{ISL} mesh, ``Careful \abbr{COTS}'' bus, and \abbr{WBS} allocation. \\ \hline
    \textbf{Clause \texttt{<5>}} \newline Life profile description & \textbf{Main Content --- Sec.\ 5} \newline Life Profile Description & Chronological lifecycle situations: on-ground \abbr{AIT}, launch, \abbr{LEOP}, commissioning, nominal modes, safe mode, and \abbr{EOL} disposal. \\ \hline
    \textbf{Clause \texttt{<6>}} \newline Environment \& constraints & \textbf{Main Content --- Sec.\ 6} \newline Environment \& Constraints & Space radiation (1,000~km vs.\ \abbr{GEO}), eclipse/cryogenic thermal regime, vacuum/outgassing, \abbr{MMOD}, mechanical loads, and regulations. \\ \hline
    \textbf{Clause \texttt{<7>}} \newline Requirements & \textbf{Main Content --- Sec.\ 7} \newline Technical Requirements & Formal verifiable requirements (\texttt{SSEG-REQ-...}) with performance tolerances, plus Traceability (\textbf{Annex A}) and Verification (\textbf{Annex B}). \\ \hline
\end{tabularx}}
\end{table}

\subsection{Normative Requirement Conventions and Identification Scheme}

In accordance with \textbf{ECSS-S-ST-00C} and \textbf{ECSS-E-ST-10-06C} Clauses 7.2.7 and 8.3 (\textbf{AD-03}), strict distinction is maintained between descriptive/informative text (Sections 1 through 6) and normative binding provisions (Section 7):
\begin{itemize}
    \item \textbf{\texttt{shall}:} Expresses a mandatory technical requirement that is contractually and technically binding on the Space Segment design and subject to formal verification per \textbf{ECSS-E-ST-10-02C} (\textbf{AD-04}).
    \item \textbf{\texttt{should}:} Expresses a non-mandatory recommendation or engineering design target where alternative solutions may be justified in the Design Justification File (\texttt{EWAS-HSB-SSEG-DDJF-03}).
    \item \textbf{\texttt{may}:} Expresses a permissible practice or optional allowance within the specified boundaries.
    \item \textbf{\texttt{can}:} Expresses a physical possibility or inherent capability of an element.
    \item \textbf{\texttt{will}} (or indicative present tense): Expresses a descriptive statement of fact, background context, or external interface condition not subject to Space Segment verification.
\end{itemize}

Furthermore, in compliance with \textbf{ECSS-E-ST-10-06C} Clause 8.3.3, ambiguous, unverifiable, or open-ended terms (such as \emph{``and/or''}, \emph{``etc.''}, \emph{``optimize''}, \emph{``minimize''}, \emph{``maximize''}, \emph{``rapid''}, \emph{``sufficient''}, \emph{``adequate''}, or \emph{``state of the art''}) are strictly excluded from the normative requirement statements in Section 7. Every technical requirement in Section 7 is assigned a unique, persistent identifier following the taxonomy:
\[
\texttt{SSEG-REQ-<CAT>-<NNN>}
\]
where \texttt{SSEG} denotes Configuration Item \texttt{CI-03} (Space Segment), \texttt{REQ} denotes a normative technical requirement, \texttt{<NNN>} is a sequential three-digit number (\texttt{001}, \texttt{002}, \dots), and \texttt{<CAT>} identifies the functional or ECSS requirement category summarized below:
\begin{itemize}
    \item \textbf{\texttt{FUN} / \texttt{MIS} / \texttt{PERF}:} Functional (\texttt{FUN}), Mission \& Constellation Geometry (\texttt{MIS}), and System Performance \& Latency (\texttt{PERF}) requirements.
    \item \textbf{\texttt{PLD} / \texttt{ISL} / \texttt{ADCS} / \texttt{OBDH}:} Infrared Sensor Payload (\texttt{PLD}), Communications \& Mandatory Inter-Satellite Links (\texttt{ISL}), Attitude \& Agile \abbr{LOS} Steering (\texttt{ADCS}), and On-Board Data Handling (\texttt{OBDH}) requirements.
    \item \textbf{\texttt{PLAT} (\texttt{STR} / \texttt{TCS} / \texttt{EPS} / \texttt{PROP}):} Integrated Platform Subsystem requirements encompassing Structures \& Mechanisms (\texttt{STR/MEC}), Thermal Control (\texttt{TCS}), Electrical Power (\texttt{EPS}), and Propulsion (\texttt{PROP}).
    \item \textbf{\texttt{IF} / \texttt{ENV} / \texttt{OPS} / \texttt{PA} / \texttt{VER}:} External Interfaces (\texttt{IF}), Space \& Launch Environment (\texttt{ENV}), Operational Modes \& Autonomy (\texttt{OPS}), Product Assurance \& ``Careful \abbr{COTS}'' (\texttt{PA}), and Verification (\texttt{VER}) requirements.
\end{itemize}
